\section{The Huntington-Hill Precedent}
\label{sec:huntington}

Congressional apportionment was historically contentious, with Congress changing formulas after nearly every census between 1790 and 1940 to favor particular states or parties \cite{balinski1982fair}. This instability ended in 1941 when Congress adopted the Huntington-Hill method \cite{huntington1928method}, which has operated unchanged for over 80 years due to its mathematical objectivity and fairness principles that transcend partisan interest.

\subsection{The Huntington-Hill Method}

The method assigns seats iteratively using priority values $P_n = \text{Population}/\sqrt{n(n-1)}$ for each state's $n$-th seat \cite{census2020apportionment}. The geometric mean denominator minimizes relative representation differences between individuals across states.

\subsection{Principles of Mathematical Fairness}

The method's success stems from six principles: \textbf{(1) Objectivity}—the formula operates deterministically on census data without human judgment; \textbf{(2) Transparency}—the mathematical foundation is fully documented and verifiable; \textbf{(3) Uniformity}—all states are treated identically; \textbf{(4) Mathematical Optimality}—it minimizes a well-defined fairness criterion; \textbf{(5) Reproducibility}—identical inputs produce identical outputs; \textbf{(6) Non-Manipulability}—the method cannot be adjusted to achieve predetermined outcomes.

\subsection{Extending to Redistricting}

Congressional redistricting faces similar challenges to pre-1941 apportionment: political manipulation, arbitrary methods, and credible claims of unfairness. Traditional criteria—population equality \cite{reynolds1964}, contiguity, and compactness \cite{polsby1991third}—can be operationalized mathematically through recursive graph bisection.

Like Huntington-Hill, recursive bisection operates deterministically on census data (tract populations and adjacencies), produces reproducible results, treats all geographic units uniformly, and optimizes a well-defined objective (edge-cut minimization). Most critically, it embodies the same principle that made Huntington-Hill successful: \textit{removing human judgment from politically sensitive decisions}. The algorithm sees only population and geography, not voting patterns or partisan affiliation, providing structural immunity to manipulation.
